\bookchapter{פרק 21: מבט מעמיק ברשתות ניידות ואבטחת מידע}

\chapterintro{אנחנו גולשים באינטרנט, משלמים ומצ'וטצ'ים כל יום בטלפון, ונראה שזה ממש נוח. אבל האם חשבת פעם, כשהיסטוריית הצ'אט, סיסמאות התשלום, התמונות שלך... כשהמידע הזה “רץ עירום” ברשת, כמה פעמים זה קורה? בכנות, רוב האנשים בכלל לא חשבו על זה, אבל ברגע שקורה משהו — כמו פריצה לחשבון או גניבה מכרטיס הבנק — אז כבר מאוחר מדי להתחרט. בפרק זה נדבר על איך להפוך את החיים הדיגיטליים שלך לפחות “שקופים”.}

\booksection{אחד, VPN: “גלימת ההיעלמות” של הרשת}

: הטלפון שלך רשת ציבורית אתר היעד. הנתבים והמתגים שבאמצע, כל שלב יכול להיות מקום שבו מישהו “מציץ” בנתונים שאתה שולח. במיוחד ב-Wi-Fi ציבורי בבתי קפה ובשדות תעופה, בכנות, זה כמעט “ריצה עירומה”.

\begin{itemize}
  \item \textbf{שימוש ב-VPN}: הטלפון שלך מנהרת VPN מוצפנת שרת VPN אתר היעד. הנתונים מוצפנים וארוזים, "מציצנים" רואים רק ג'יבריש. בנוסף, אתר היעד רואה את כתובת ה-IP של שרת ה-VPN ולא את ה-IP האמיתי שלך.
  \item \textbf{אז מה VPN באמת עושה? בקיצור: ראשון,}מגן על הפרטיות שלך ב-Wi-Fi ציבורי
\end{itemize}

; שניים,\textbf{עוקף מגבלות גיאוגרפיות}(למשל צפייה בתוכן שזמין רק במדינות מסוימות); שלישי,\textbf{מונע מספק האינטרנט (ISP) שלך לתעד את הרגלי הגלישה שלך}. אבל אני חייב להזהיר:\textbf{אל תשתמש ב-VPN חינמי!}“חינמי” לרוב אומר שהנתונים שלך הם ה“סחורה” שלהם. יש לי חבר שהשתמש ב-VPN חינמי, ואז הטלפון שלו התחיל לקפוץ עם פרסומות, ואחר כך התברר שה-VPN מכר את נתוני הגלישה שלו למפרסמים. לכן בבחירת VPN, או שתבחר במותג גדול בתשלום (כמו NordVPN, ExpressVPN), או פשוט אל תשתמש.\textbf{בנוסף, VPN אינו תרופת פלא. הוא מגן רק על נתונים “במעבר”, אם אתר היעד עצמו אינו מאובטח, או שכבר יש תוכנה זדונית בטלפון שלך, VPN לא יעזור לך.}שניים, “מלכודת הדבש” של Wi-Fi ציבורי — מסוכן יותר ממה שאתה חושב

אולי נראה לך, מה כבר יכול לקרות מחיבור ל-Wi-Fi חינמי? אספר לך סיפור אמיתי: חוקרי אבטחה ערכו ניסוי, הציבו בשדה תעופה נקודת חיבור בשם “Starbucks\_Free”, ובפחות משעה התחברו מעל מאה אנשים. למעשה, נקודת החיבור הזו הוקמה על ידי החוקרים עצמם, וכל

\booksection{שניים, “מלכודת הדבש” של Wi-Fi ציבורי — מסוכן יותר ממה שאתה חושב}

שעברה דרכה — למשל תוכן של דפי HTTP שביקרת בהם, תקשורת App ללא TLS — הם באמת יכלו ליירט ולצפות בה.\textbf{אבל שים לב להנחת יסוד קריטית: היום WeChat, בנקאות אונליין, Alipay שאתה משתמש בהם יומיומית, רובם כבר עוברים הצפנת HTTPS. במילים אחרות, גם אם התוקף מיירט את המנות, הוא רואה רק ג'יבריש,}ולא יכול לקרוא ישירות את תוכן הצ'אט או סיסמת התשלום שלך

. עם זאת, הם עדיין יכולים לדעת לאילו אתרים נכנסת (שם הדומיין), ואם App או דף ישן עדיין משתמשים בתעבורת HTTP גלויה, התוכן יהיה חשוף לחלוטין. התקפה זו נקראת “התקפת איש באמצע” (Man-in-the-Middle), והעיקרון הוא שהתוקף “מכניס יד” בינך לבין הנתב האמיתי. מפחיד יותר, הטלפון של אנשים רבים מוגדר ל“חיבור אוטומטי ל-Wi-Fi מוכר”, והתוקף רק צריך לקרוא לנקודה שלו בשם Wi-Fi שהתחברת אליו בעבר (כמו “ChinaNet” או “CMCC”), והטלפון שלך יתחבר אליו אוטומטית.\textbf{אז איך מתגוננים? כמה עצות מעשיות:}סמוך רק על HTTPS

: תחילת שורת הכתובת היא https://, ה-“s” הנוסף מסמן Secure. בעת גלישה בכל אתר שדורש הזנת סיסמה, אם אתה רואה http:// (ללא s), עצור מיד ואל תזין! היום דפדפנים מרכזיים מציגים אזהרת “לא מאובטח” לפני אתרים שאינם HTTPS, כשרואים סימון זה יש להיזהר.

\begin{itemize}
  \item \textbf{כבה "חיבור אוטומטי"}: בין אם iPhone ובין אם Android, בהגדרות Wi-Fi כבה את "הצטרפות אוטומטית" או "חיבור אוטומטי". עדיף ללחוץ ידנית בכל פעם מאשר לתת לטלפון להחליט.
  \item \textbf{השתמש בחבילת הגלישה של הטלפון לפעולות רגישות}: העברות כספים, התחברות לבנקאות אונליין, הזנת מספר תעודת זהות... את הפעולות האלה בצע ב-4G/5G שלך, אל תחסוך בתעבורה.
  \item \textbf{"שכח" Wi-Fi שאינו בשימוש תכוף}: כשחוזרים מנסיעת עסקים או טיול, "שכח" את רשתות ה-Wi-Fi של בתי המלון ושדות התעופה שהתחברת אליהן זמנית. כך בפעם הבאה הטלפון לא יתחבר אוטומטית לנקודה מזויפת באותו שם.
  \item \textbf{שלוש, כיצד HTTPS מגן עליך? — לחיצת יד SSL/TLS}הזכרנו קודם את HTTPS, אז איך הוא באמת עובד? בכל פעם שאתה נכנס לאתר HTTPS, הדפדפן והשרת מבצעים "לחיצת יד SSL/TLS", כל התהליך נראה בערך כך:
\end{itemize}

\booksection{שלוש, כיצד HTTPS מגן עליך? — לחיצת יד SSL/TLS}

השרת עונה: “בסדר, בחרתי מהרשימה שלך את סכמת ההצפנה הזו, הנה התעודה הדיגיטלית שלי, שמכילה את המפתח הציבורי שלי, חתומה ומאומתת על ידי רשות CA כלשהי.”

\begin{enumerate}
  \item הדפדפן מאמת את התעודה: האם היא הונפקה על ידי CA מהימן? האם פג תוקפה? האם שם הדומיין תואם?
  \item לאחר אימות, שני הצדדים משתמשים במפתח הציבורי כדי להצפין ולהסכים על "מפתח סשן", וכל התקשורת לאחר מכן מוצפנת במפתח זה.
  \item זה כמו שאתה הולך לבנק: אתה מגיש את תעודת הזהות שלך (תעודה), הפקיד מאמת את זהותך, ואז אתם מסכימים על סיסמה, וכל השיחה לאחר מכן רק אתם יכולים להבין. גם אם מישהו מצותת בצד, הוא שומע רק ג'יבריש מוצפן.
  \item אבל שים לב: HTTPS מבטיח רק ש"תהליך ההעברה" מוצפן, לא שאתר היעד הוא איש טוב. גם אתרי פישינג יכולים להשתמש ב-HTTPS — סמל המנעול הקטן רק מעיד שהתחברת ל"אתר כלשהו", לא שהאתר הזה הוא זה שרצית. לכן עדיין צריך לבדוק בעצמך שהכתובת נכונה.
\end{enumerate}

ארבע, אבטחת DNS: האם “ספר הטלפונים של הרשת” שלך מאובטח?

DNS (מערכת שמות דומיינים) היא כמו ספר הטלפונים של האינטרנט: אתה מזין “baidu.com”, ו-DNS עוזר לך למצוא לאיזו כתובת IP הוא מתאים (למשל 110.242.68.66), ואז הדפדפן יכול להתחבר. הבעיה היא ששאילתות DNS מסורתיות הן ב

\booksection{ארבע, אבטחת DNS: האם “ספר הטלפונים של הרשת” שלך מאובטח?}

,\textbf{לא מוצפנות}! כלומר, ספק האינטרנט שלך, ספק ה-Wi-Fi, או כל נקודת ביניים, יכולים לראות לאילו אתרים נכנסת.\textbf{גרוע יותר זה “חטיפת DNS”: התוקף משנה את שלב ה-DNS, ומפנה את “baidu.com” לאתר פישינג מזויף. אתה מזין את הכתובת הנכונה, אבל מגיע לדף התחברות זהה לחלוטין אך מזויף, וברגע שתזין את שם המשתמש והסיסמה — נגמר.}למרבה המזל, כיום יש פתרונות DNS מוצפנים:

DNS over HTTPS (DoH)

: מסתיר שאילתות DNS בתוך תעבורת HTTPS, מעורבב בתעבורת גלישת אינטרנט רגילה, כך שלא ניתן להבחין אילו הן שאילתות DNS.

\begin{itemize}
  \item \textbf{DNS over TLS (DoT)}: משתמש בפרוטוקול TLS להצפנה ייעודית של שאילתות DNS, פורט עצמאי, מנהל הרשת יכול לדעת שהשתמשת ב-DoT אך לא רואה את התוכן.
  \item \textbf{אני אישית ממליץ להשתמש ב-1.1.1.1 של Cloudflare או ב-8.8.8.8 של Google כשרתי DNS, שניהם תומכים ב-DNS מוצפן. בהגדרות "DNS פרטי" של Android פשוט הזן "one.one.one.one" או "dns.google", קל מאוד. ב-iPhone זה קצת יותר מסובך, צריך להתקין פרופיל תיאור או להגדיר דרך App, אבל שווה לעשות את זה.}חמישה, האם "גלישה במצב חסוי" באמת מסתירה אותך?
\end{itemize}

קודם כל, המסקנה:

\booksection{חמישה, האם “גלישה במצב חסוי” באמת מסתירה אותך?}

. הם נמצאים בין מוצר נישה לחובבים למוצר צריכה המוני, ומתאימים לקהל מסוים.\textbf{מה הוא לא יכול לעשות:}

האתר שאתה מבקר

עדיין רואה שהגעת

\begin{itemize}
  \item , כתובת ה-IP, פרטי המכשיר והתנהגות הגלישה שלך ברורים לאתר.\textbf{ספק האינטרנט שלך}עדיין רואה לאילו אתרים נכנסת
  \item , רק לא רואה את התוכן המדויק (אם משתמשים ב-HTTPS).\textbf{מנהל הרשת של החברה/בית הספר}יכול גם לנטר
  \item את התנהגות הגלישה שלך.\textbf{קבצים שהורדת, סימניות שהוספת}יישמרו
  \item , ולא ייעלמו עם סגירת החלון.\textbf{לכן זכור: מצב חסוי הוא "מוחק מקומי", לא "בגד היעלמות ברשת". אם באמת צריך להסתיר עקבות גלישה, צריך שילוב של VPN + דפדפן Tor, אבל לשימוש יומיומי, הידיעה של הדברים האלה מספיקה.}שישה, פרטיות בדפדפן: אל תיתן לדפדפן למכור אותך
\end{itemize}

אולי לא מודע לזה, אבל הדפדפן עצמו הוא אוסף מידע ענק. עוגיות, היסטוריית גלישה, מטמון, השלמה אוטומטית, הרשאות מיקום... כל אחד יכול להדליף את הפרטיות שלך. כמה עצות:

\booksection{שישה, פרטיות בדפדפן: אל תיתן לדפדפן למכור אותך}

: מדי פעם, כנס להגדרות ונקה עוגיות, מטמון והיסטוריה. לא כדי למנוע “מעקב” — זה כבר לא עובד עם ניקוי עוגיות — אלא כדי למנוע מהטלפון להיות איטי. מטמון הדפדפן יכול לתפוס מספר ג'יגה-בייט.

\begin{itemize}
  \item \textbf{כבה את האשליה של "אל תעקוב"}: האפשרות "שלח בקשת 'אל תעקוב'" בדפדפן היא בעצם "בקשה" מהאתר לא לעקוב אחריך. אבל אם האתר מקשיב או לא זה עניין אחר לחלוטין, רוב האתרים פשוט מתעלמים. אל תסמוך על זה.
  \item \textbf{השתמש בדפדפנים המגנים על פרטיות}: Firefox Focus, דפדפן DuckDuckGo, Brave — דפדפנים אלו חוסמים עוקבים ופרסומות כברירת מחדל, הרבה יותר נקיים מההגדרות ברירת המחדל של Chrome ו-Safari.
  \item \textbf{גם מנוע החיפוש חשוב}: השתמש ב-DuckDuckGo במקום Baidu/Google לחיפוש, הוא לא רושם את היסטוריית החיפוש שלך ולא בונה פרופיל משתמש. בכנות, תוצאות החיפוש פחות מדויקות מ-Google, אבל מספיקות לחיפוש יומיומי.
  \item \textbf{בדוק הרשאות App}: כנס להגדרות הטלפון ובדוק אילו הרשאות ניתנו לדפדפן שלך. האם הוא באמת צריך גישה למיקום, למיקרופון, לאנשי הקשר שלך? כנראה שלא.
  \item \textbf{שבעה, רשימת אבטחה מעשית: התחל להגן על עצמך עכשיו}אמרתי הרבה, אז הנה רשימת אבטחה פשוטה, חצי שעה וסיימת, ואבטחת הטלפון שלך תשתפר משמעותית:
\end{itemize}

\booksection{שבעה, רשימת אבטחה מעשית: התחל להגן על עצמך עכשיו}

: iPhone בהגדרות - Wi-Fi, Android בהגדרות Wi-Fi.

\begin{enumerate}
  \item \textbf{החלף ל-DNS מוצפן}: Android ישירות הגדר DNS פרטי ל-"dns.google", iPhone באמצעות App או פרופיל תיאור להגדרת DoH של 1.1.1.1.
  \item \textbf{הפעל אימות דו-שלבי (2FA)}: WeChat, Alipay, Apple ID, חשבון Google, כל מה שאפשר להפעיל 2FA — הפעל. קוד SMS הוא כבר ברמה הנמוכה ביותר של 2FA, עדיף להשתמש ב-TOTP (כמו Google Authenticator, Microsoft Authenticator).
  \item \textbf{בדוק הרשאות App}: עבור על ההרשאות של כל App, במיוחד מיקום, אנשי קשר, SMS, מיקרופון, מצלמה. תגלה שאפליקציות פנס מסוימות דורשות את אנשי הקשר שלך — הסר מיד.
  \item \textbf{נקה נתוני דפדפן תקופתית}: קבע תזכורת ביומן, פעם בחודש, נקה מטמון ועוגיות.
  \item \textbf{אל תשתמש בכבלי נתיבת טעינה ציבורית}: יציאת USB יכולה גם לטעון וגם להעביר נתונים, עמדת טעינה זדונית (נקראת "Juice Jacking") יכולה לגנוב נתונים בזמן שאתה טוען. קח עם ך מטען וכבל משלך.
  \item \textbf{עשה זאת עכשיו: בצע "בדיקת אבטחה" לטלפון שלך}שדרוג DNS
\end{enumerate}

\begin{tcbpractice}{: Android: הגדרות > חיבורים ושיתוף > DNS פרטי > הזן "one.one.one.one" או "dns.google". iPhone: חפש ב-Safari "1.1.1.1 DNS profile" והתקן פרופיל תיאור, או הורד את אפליקציית 1.1.1.1 של Cloudflare.}
\begin{enumerate}
  \item \textbf{ביקורת Wi-Fi}: פתח את רשימת ה-Wi-Fi, "שכח" כל נקודות החיבור שלא מזוהות, שאינן בשימוש תכוף, או בעלות שם מוזר.
  \item \textbf{ניקוי הרשאות יסודי}: כנס להגדרות > פרטיות/ניהול הרשאות, עבור על הרשאות כל App. סמן App שדרישות ההרשאות שלהן אינן סבירות (למשל מחשבון שדורש אנשי קשר), שקול להסיר או להשתמש בתחליף.
  \item \textbf{בדוק אבטחת חשבון}: פתח את דף האבטחה של Apple ID או חשבון Google, בדוק אם יש התחברויות ממכשירים לא מזוהים, ודא ש-2FA מופעל, בדוק שאנשי הקשר לשחזור אמינים.
  \item \textbf{פרק 22: כיצד AI עובד בטלפון}ברוך הבא ל"מופע הקסמים" הכי מגניב של הטלפון, הכוכב הוא: בינה מלאכותית. התמונות היפות ש"נלחצות במקלדת אחת", העוזר הקולי ש"פותח את הפה", שיטת הקלט ש"מבינה אותך יותר ככל שמשתמשים יותר" — מאחורי כולם עומד המוח הנסתר הזה.
\end{enumerate}

\end{tcbpractice}
